% Generated draft: requires mathematical and source review.
\documentclass[11pt]{article}
\usepackage[margin=25mm]{geometry}
\usepackage{fontspec}
\setmainfont{Latin Modern Roman}
\usepackage{amsmath,amssymb,mathtools,unicode-math}
\setmathfont{Latin Modern Math}
\usepackage{enumitem,needspace,xurl,hyperref}
\setlength{\parindent}{0pt}
\setlength{\parskip}{5pt}
\setlength{\emergencystretch}{3em}
\newcommand{\N}{\mathbb{N}}
\newcommand{\Z}{\mathbb{Z}}
\newcommand{\Q}{\mathbb{Q}}
\newcommand{\R}{\mathbb{R}}
\newcommand{\C}{\mathbb{C}}
\newcommand{\F}{\mathbb{F}}
\newcommand{\A}{\mathbb{A}}
\newcommand{\PP}{\mathbb{P}}
\newcommand{\E}{\mathbb{E}}
\newcommand{\eps}{\varepsilon}
\newcommand{\dd}{\,\mathrm{d}}
\begin{document}
\section*{UnsolvedMath ID: 30000102}
\textbf{Source problem code:} catalogue entry
\subsection{Definitions and mathematical statement}
The source asks whether every complex Banach space \(X\) satisfying
\[
 \Lambda_{\mathrm{Bloch},\ell^1}(X)=\ell^{(2,1)}
\]
must be isomorphic to a Hilbert space. The source does not define the invariant on the left, the sequence-space convention on the right, or the meaning of equality. In particular, the notation \(\ell^{(2,1)}\) must not be assigned a Lorentz or mixed/block norm merely from its typography. The title ``Bloch Spectra'' does not supply the missing construction.

The conclusion has an unambiguous standard interpretation: there should exist a complex Hilbert space \(H\) and a bounded complex-linear bijection \(T:X\to H\) with bounded inverse. Equivalently, \(X\) should admit a positive Hermitian inner product \([\cdot,\cdot]\) and constants \(0<c\leq C<\infty\) such that
\[
 c\|x\|_X^2\leq[x,x]\leq C\|x\|_X^2\qquad(x\in X).
\]
This asks for an equivalent Hilbertian norm, not that the given norm already satisfy an inner-product identity.

Thus the intended universal implication is recorded with its exact displayed hypothesis, while its mathematical verification requires the original definitions, index/weight conventions, and whether equality means equality of sequence sets, normed identification, or norm equivalence. No substitute for \(\Lambda_{\mathrm{Bloch},\ell^1}(X)\) is invented.

The supplied literature triage dated 21 August 2026 discusses sufficiency for \(L^2\), counterexamples among symmetric spaces with trivial Boyd indices, and an incomplete classification of rearrangement-invariant spaces. It does not explain how that classification relates to the displayed Bloch-invariant implication for arbitrary complex Banach spaces. The reported partial results are therefore retained as a possible curation-context mismatch, not asserted as theorems about this undefined hypothesis. The supplied status is partially solved, with the exact scope unresolved.
\subsection{Short English statement}
The target is an equivalent Hilbert norm for a complex Banach space satisfying the displayed Bloch sequence-space identity. Both sides of that hypothesis lack definitions in the record, and the supplied literature assessment discusses a different apparent classification framework. The gap is explicit rather than filled by guessing the notation.
\subsection{Sources}
\begin{itemize}
\item Local UnsolvedMath registry, permanent ID \(30000102\), exact source code \texttt{catalogue entry}, Oberwolfach Reports Open Problems. Statement status \texttt{corrected\_verified}; verification only removes a question label and balances the TeX, without supplying invariant definitions.
\item ``Spectral Theory in Banach Spaces and Harmonic Analysis (2005),'' Oberwolfach Reports, pp. 1883--1970, DOI \texttt{10.4171/owr/2004/36}: \url{https://doi.org/10.4171/owr/2004/36}. Proposed year and DOI refer to 2004, while the citation label says 2005.
\item Supplied ``de Pagter--Sukochev'': \url{https://doi.org/10.1016/S0022123603003860}. The other supplied source is the original report DOI above. No full invariant definition or explicit connecting theorem is included.
\item Literature checked 21 August 2026; supplied status \texttt{partially\_solved}, with an assessment about rearrangement-invariant spaces and Boyd indices. Registry created 25 August 2026 and updated 16 September 2026. This dated triage requires independent verification before citation.
\item Registry curation attribution: Ulam AI / UnsolvedMath Contributors, CC BY 4.0. Underlying report and literature retain their supplied provenance; this statement-only draft requires human review.
\end{itemize}
\end{document}
